\thispagestyle{plain}
\begin{tikzpicture}[remember picture,overlay]
  \fill[ink] (current page.north west) rectangle ([yshift=-8.6cm]current page.north east);
  \fill[teal] ([yshift=-8.6cm]current page.north west) rectangle ([yshift=-8.68cm]current page.north east);
  \node[anchor=north west,inner sep=0pt] at ([xshift=2.5cm,yshift=-1.9cm]current page.north west)
    {\includegraphics[height=1.15cm]{logo/nonos-logo-cyan.pdf}};
  \node[anchor=north west,text width=15.4cm,inner sep=0pt,text=white] at
    ([xshift=2.5cm,yshift=-3.9cm]current page.north west)
    {\fontsize{25}{29}\selectfont\bfseries A computer that can prove\\what it runs\par
     \vspace{0.9em}
     {\fontsize{12.5}{16}\selectfont\color{cyan}\mdseries Release 0.9.3: proof for the software we ship,\\and a box for everyone else's\par}
     \vspace{1.1em}
     {\footnotesize\ttfamily\color{white!70!ink} whitepaper on release 0.9.3 \quad repository commit \quad 30 October 2026\par}};
\end{tikzpicture}
\vspace*{6.3cm}

{\footnotesize\color{mute}\ttfamily https://nonos.software \quad https://github.com/NON-OS/microkernel-preview \quad https://verify.nonos.software}

\vspace{1.1em}
{\large\bfseries\color{teal} Abstract}\par\vspace{0.2em}
\noindent\nonos{} is an operating system for x86\_64 computers that tries to make its software state checkable, by the machine and by someone else, before that software is trusted with keys, money or traffic. A Rust capability microkernel keeps page tables, scheduling, IPC and 183 checked system calls in ring 0. Drivers, the network stack, the desktop and a wallet run in ring 3 as 152 signed capsules, each admitted at spawn by two signatures and a STARK membership proof under one policy root.

Release 0.9.3 adds what 0.9.2 left open: programs the project did not write. A Linux program opens its own window on the desktop, inside a box named by its signed Marketplace listing: a window, its own home folder, the chosen anonymity network, a terminal, the right to start programs. The program holds no capability of its own. The kernel sends every call it makes to one supervisor capsule, which applies the box, and keeps its own count of the answers. We show what the system refused when single hostile attempts were made from inside a program, and what it did not, and list the three known ways around the box.

The paper says what has been proved, what has been booted, and what is only designed. No independent audit has been done.


\begin{callout}
\textbf{Reading this paper.} Every technical statement cites a file in the repository at the commit above, with line numbers where they help, in a numbered footnote. A claim we could not tie to the tree is not made. Where two documents in the tree disagree, we say so and follow the code. Section~\ref{sec:why} states the problem. Sections~\ref{sec:arch} to~\ref{sec:capsules} describe the kernel and how it admits software. Section~\ref{sec:box} is the centre of this release: running programs we did not write. Section~\ref{sec:status} says what has run and what has not. Nothing in this paper is an audit, and none has been done.
\end{callout}
\clearpage
